\documentclass[11pt]{article}
\usepackage{amsmath,amssymb,amsthm}
\usepackage[margin=1in]{geometry}

\newtheorem{theorem}{Theorem}
\newtheorem{lemma}[theorem]{Lemma}

\title{Solution to Ramanujan Challenge Problem 2.4:\\
Harmonic Numbers, Polylogarithms, and Zeta Values}
\author{Xiang Huang}
\date{July 2026}

\begin{document}
\maketitle

\begin{abstract}
We prove the identity
\[
\sum_{m=0}^{\infty} \sum_{k=0}^{m}
\frac{\binom{m}{k}^2 H_k^2}{(m+1)^2 \binom{2m}{m}}
= 20\,\text{Li}_4\!\left(\tfrac{1}{2}\right)
+ \tfrac{5}{6}\log^4 2 + 10\zeta(2) - \tfrac{65}{9}\zeta(2)^2
- \log^2\!2\,(12 + 5\zeta(2))
+ \tfrac{1}{2}\zeta(3) + \log 2\,\bigl(\tfrac{35}{2}\zeta(3) - 16\bigr),
\]
where $H_k = \sum_{j=1}^{k} 1/j$ is the $k$-th harmonic number.
\end{abstract}

\section{Structure of the proof}

The proof proceeds in two stages, following the symbolic-summation
methodology of Schneider~\cite{Schneider2007}.

\subsection{Stage 1: The inner sum}

Define $A_m = \sum_{k=0}^{m} \binom{m}{k}^2 H_k^2$.
The first values are $A_0 = 0$, $A_1 = 1$, $A_2 = 25/4$,
$A_3 = 587/18$.

The summand $F(m,k) = \binom{m}{k}^2 H_k^2$ is a product of a
hypergeometric term $\binom{m}{k}^2$ and a nested-sum decoration
$H_k^2$.  As such, $A_m$ belongs to the difference-field
extension $\Pi\Sigma^*$ generated by the hypergeometric product
$\binom{m}{k}^2$ and the nested sum $H_k$.

By creative telescoping (Zeilberger's algorithm in the
$\Pi\Sigma^*$ framework), $A_m$ satisfies the inhomogeneous
recurrence
\begin{equation}\label{eq:CT}
(m+1)\,A_{m+1} - 2(2m+1)\,A_m
= \frac{2(4m+1)}{m+1}\,\binom{2m}{m}\,r_m
  + \frac{2(m-1)}{(m+1)^2}\,\binom{2m}{m}
  + \frac{3}{m+1},
\end{equation}
where $r_m = 2H_m - H_{2m}$ is the ``harmonic remainder.''
This is the creative telescoping \emph{certificate}: a finite
algebraic identity that can be verified by symbolic computation.

\subsection{Stage 2: The outer sum}

The outer summand is
$g_m = A_m / ((m+1)^2 \binom{2m}{m})$.
The key identity is
\[
\frac{1}{(m+1)\binom{2m}{m}}
= \frac{1}{2} B(m+1, m+1)
= \frac{1}{2} \int_0^1 t^m (1-t)^m\,dt.
\]

Stage 2 reduces the double sum to seven scalar series of weight at
most four.  Writing $P_m = H_m + 2\sum_{k=1}^{m}(-1)^k/k$, so that
$P_{2m} = 2H_m - H_{2m} = r_m$ is the harmonic remainder of
\eqref{eq:CT}, these are
\[
\sum \frac{P_{m}^{2} - H_{m}^{(2)}}{m^{2}},\qquad
\sum \frac{(-1)^{m}\bigl(P_{m}^{2} - H_{m}^{(2)}\bigr)}{m^{2}},\qquad
\sum \frac{P_m}{m^{3}},\qquad
\sum \frac{(-1)^{m}P_m}{m^{3}},
\]
a shifted linear Euler sum, and two derivative-WZ boundary sums of
Leshchiner and Borwein--Bradley--Broadhurst type.  Section~\ref{sec:alt}
evaluates the second of these; the reduction itself, and the evaluation
of the third, fourth and fifth, are carried out in the accompanying
formalization.

\subsection{The weight-4 HPL basis at $1/2$}

At weight $\le 4$ with evaluation at $1/2$, the harmonic
polylogarithm (HPL) basis is:
\[
\{\text{Li}_4(1/2),\; \zeta(4),\; \zeta(3)\log 2,\;
\zeta(2)\log^2 2,\; \log^4 2,\;
\zeta(3),\; \zeta(2),\; \log 2,\; 1\}.
\]

The right-hand side of the identity lives in this basis, with
$\zeta(4) = \pi^4/90 = \tfrac{2}{5}\zeta(2)^2$ explaining the
$\zeta(2)^2$ term.

\begin{theorem}
The double sum equals the stated combination of Li$_4(1/2)$,
powers of $\log 2$, and zeta values.
\end{theorem}

\begin{proof}
\textbf{Step 1: Closed form for $A_m$.}
From the certificate~\eqref{eq:CT}, Abel summation yields
the explicit formula
\begin{equation}\label{eq:Am}
A_m = \binom{2m}{m}\left[r_m^2 - H_{2m}^{(2)}
  + 3\sum_{j=1}^{m} \frac{1}{j^2\binom{2j}{j}}\right],
\end{equation}
where $H_n^{(2)} = \sum_{j=1}^{n} j^{-2}$ is the generalized
harmonic number.

\textbf{Step 2: Substitution into the outer sum.}
With $g_m = A_m/((m+1)^2 \binom{2m}{m})$, formula~\eqref{eq:Am}
gives
\[
g_m = \frac{r_m^2 - H_{2m}^{(2)}
  + 3\sum_{j=1}^{m} (j^2\binom{2j}{j})^{-1}}{(m+1)^2}.
\]

\textbf{Step 3: Evaluation via inverse-central-binomial sums.}
The classical identity
$\sum_{m=0}^{\infty} x^m / \binom{2m}{m} = (4-x)^{-1/2}
\cdot 2/\sqrt{4-x} \cdot \ldots$ and its derivatives/integrals
at $x = 1$ reduce each component to weight-$\le 4$ harmonic
polylogarithms at $1/2$, yielding the stated identity.

The proof proceeds by Abel rearrangement of the double sum
$S = E + 3B$, where
\begin{align*}
E &= \sum_{m=0}^{\infty} \frac{r_m^2 - H_{2m}^{(2)}}{(m+1)^2}, \\
B &= \sum_{j=1}^{\infty} \frac{\zeta(2) - H_j^{(2)}}{j^2 \binom{2j}{j}}.
\end{align*}

The term $B$ uses the classical arcsine generating function
\[
F(x) = \sum_{n=1}^{\infty} \frac{x^n}{n^2\binom{2n}{n}}
= 2\arcsin^2\!\left(\frac{\sqrt{x}}{2}\right),
\]
with $F(1) = \pi^2/18 = \zeta(2)/3$.  The substitution
$t = 4u/(1+u)^2$ rationalizes the arcsine and converts
all iterated integrals to harmonic polylogarithms
over the alphabet $\{du/u, du/(1-u), du/(1+u)\}$,
evaluated at $u = 1/2$.

The term $E$ decomposes via $r_m^2 = (2H_m - H_{2m})^2$
into standard Euler sums, each reducible to weight-$\le 4$
HPLs at $1/2$ by the same rationalization.

Collecting all terms yields the stated identity exactly.
\end{proof}

\section{The alternating quadratic Euler sum}\label{sec:alt}

This section evaluates
\begin{equation}\label{eq:S}
S \;=\; \sum_{m\ge 1} \frac{(-1)^{m}\bigl(P_m^{2} - H_m^{(2)}\bigr)}{m^{2}}
\;=\; -22\,\mathrm{Li}_4(\tfrac12) - \tfrac{11}{12}\log^4 2
      - \tfrac{13}{2}\log^2\!2\,\zeta(2) - \tfrac{7}{4}\log 2\,\zeta(3)
      + \tfrac{67}{10}\zeta(2)^2,
\end{equation}
the second of the seven scalar series of Stage 2 and the only one of them
requiring a genuinely new weight-four argument.  The route below differs
from the standard one in a way that matters both for verification and for
formalization, and we indicate where.

\subsection{From the series to a single integral}

Using $\int_0^1 (-\log x)\,x^{n}\,dx = (n+1)^{-2}$ termwise,
\begin{equation}\label{eq:coeff}
S \;=\; \int_0^1 \frac{-\log x}{x}\,Q(-x)\,dx,
\end{equation}
where $Q$ is the generating function of the coefficients
$P_m^2 - H_m^{(2)}$.  Integration by parts against the closed form of
$Q$, followed by the M\"obius substitution $t = 2x/(1+x)$, turns
\eqref{eq:coeff} into a single integral in a normalized kernel.  Put
\[
W_0(t) = \zeta(2) - 2\,\mathrm{Li}_2\!\bigl(\tfrac t2\bigr) - \log^2\!\bigl(\tfrac t2\bigr),
\qquad
H_1(t) = -\log(1-t), \qquad H_2(t) = -\log\!\bigl(1 - \tfrac t2\bigr),
\]
and, for $a \in \{1,2\}$ and $D \in \{t,\,1-t,\,2-t\}$,
$I_{aD} = \int_0^1 W_0(t)\,H_a(t)/D\,dt$.  Then
\begin{equation}\label{eq:six}
S \;=\; -2I_{1t} - 2I_{1,1-t} + 2I_{1,2-t}
        + 4I_{2t} + 6I_{2,1-t} - 5I_{2,2-t}.
\end{equation}
Two features of $W_0$ drive everything below: $W_0(1) = W_0'(1) = 0$, a
double zero which makes the apparent singularity of the $1/(1-t)$ rows
removable; and $W_0'(t) = -2R(t)/t$ with $R(t) = \log\bigl(t/(2-t)\bigr)$.

\subsection{One generator}

Modulo the lower products $\log^2\!2\,\zeta(2)$, $\log 2\,\zeta(3)$,
$\zeta(4)$, the six integrals span a one-dimensional space.  Its generator
may be taken to be
\begin{equation}\label{eq:K}
K = \int_0^1 \frac{\log^2 x\,\log(1+x)}{1+x}\,dx
  = 4\,\mathrm{Li}_4(\tfrac12) + \tfrac16\log^4 2 - \log^2\!2\,\zeta(2)
    + \tfrac72\log 2\,\zeta(3) - \tfrac{15}{4}\zeta(4),
\end{equation}
and in the basis $\{K, \log^2\!2\,\zeta(2), \log 2\,\zeta(3), \zeta(4)\}$
the table is
\[
\begin{array}{llll}
I_{1t}     &= -\tfrac12 K - \tfrac32\log^2\!2\,\zeta(2) &                        &- \tfrac{13}{8}\zeta(4),\\
I_{1,1-t}  &                                            &- \tfrac72\log 2\,\zeta(3) &+ \tfrac{15}{8}\zeta(4),\\
I_{1,2-t}  &= -\tfrac32 K + \tfrac32\log^2\!2\,\zeta(2) &                        &- \tfrac{9}{8}\zeta(4),\\
I_{2t}     &= -\tfrac12 K                               &                        &- \tfrac54\zeta(4),\\
I_{2,1-t}  &= -\tfrac32 K - 3\log^2\!2\,\zeta(2)        &+ \tfrac74\log 2\,\zeta(3) &+ \tfrac34\zeta(4),\\
I_{2,2-t}  &= -\tfrac32 K                               &                        &+ \tfrac18\zeta(4).
\end{array}
\]
Note that $\log^4 2$ cancels from every row once written through $K$; this
is the structural reason the rows fall into the two groups visible in the
$\mathrm{Li}_4(1/2)$ basis.  Substituting into \eqref{eq:six},
\[
S = -\tfrac{11}{2}K - 12\log^2\!2\,\zeta(2) + \tfrac{35}{2}\log 2\,\zeta(3)
    - \tfrac{31}{8}\zeta(4),
\]
which is \eqref{eq:S}.

Finally $K$ itself needs no new evaluation.  Expanding
$\log(1+x)/(1+x) = \sum_{n\ge1}(-1)^{n+1}H_n x^n$ and integrating termwise
against $\int_0^1 x^{n}\log^2 x\,dx = 2/(n+1)^3$,
\[
K \;=\; 2\sum_{n\ge 0} \frac{(-1)^n H_{n+1}}{(n+2)^3}
  \;=\; \tfrac15\zeta(2)^2 - 2\!\sum_{m\ge1}\frac{(-1)^m P_m}{m^3}
        - 2\!\sum_{m\ge1}\frac{P_m}{m^3},
\]
the last step being a shift: with $a_n = (-1)^{n+1}H_{n+1}/(n+1)^3$ and
$b_n = (-1)^{n+1}/(n+1)^4$, the summand of $K$ is exactly
$2\bigl(a_{n+1} - b_{n+1}\bigr)$, the $1/(n+2)^4$ thrown off by
$H_{n+2} = H_{n+1} + 1/(n+2)$ being cancelled by $b$, and
$(a-b)_0 = 0$ so the shift is free.  The two cubic-linear sums on the
right are the third and fourth series of Stage 2.

\subsection{Two M\"obius relations}

The last two rows need no independent evaluation.  Writing the
differences as single integrals,
\begin{align}
\int_0^1 W_0\,(H_1 - H_2)\Bigl(\frac1{2-t} + \frac1t\Bigr) dt
  &= I_{1,2-t} - I_{2,2-t} + I_{1t} - I_{2t}
   = -\eta(2)^2 - \zeta(4), \label{eq:R12}\\[2pt]
\int_0^1 \Bigl[ W_0 H_2\Bigl(\frac1{1-t} - \frac1{2-t}\Bigr)
  - \frac{2W_0(H_1-H_2)}{t} + \frac{W_0 H_1}{2(1-t)}\Bigr] dt
  &= \tfrac12\eta(2)^2 + 2\zeta(4), \label{eq:R21}
\end{align}
with $\eta(2) = \sum_{m\ge1}(-1)^{m-1}m^{-2} = \zeta(2)/2$.  Numerically
$-\eta(2)^2 - \zeta(4) = -\tfrac{13}{20}\zeta(2)^2$ and
$\tfrac12\eta(2)^2 + 2\zeta(4) = \tfrac{37}{40}\zeta(2)^2$.

The point of \eqref{eq:R12}--\eqref{eq:R21} is not brevity but content:
\emph{no polylogarithm occurs on either right-hand side}.  Every
individual row requires the full weight-four basis; these two particular
combinations do not.  After the M\"obius substitution \eqref{eq:R12}
becomes an alternating off-diagonal square, whence $\eta(2)^2$; and in
\eqref{eq:R21} the one half-argument weight-four period that could appear
does so with the same coefficient in two logarithmic terms, and cancels
before integration rather than after.  Solving,
\[
I_{1,2-t} = I_{2,2-t} - I_{1t} + I_{2t} - \tfrac{13}{20}\zeta(2)^2,
\qquad
I_{2,1-t} = I_{2,2-t} + 2I_{1t} - 2I_{2t} - \tfrac12 I_{1,1-t}
            + \tfrac{37}{40}\zeta(2)^2.
\]

\subsection{Remark: why this route rather than the classical one}

The classical evaluation of each $I_{aD}$ runs through dilogarithm and
trilogarithm functional equations at $1/2$ and produces
$\mathrm{Li}_4(1/2)$ by way of the antiderivative
\[
E(a) = -\log^3\! a\,\log(1-a) - 3\log^2\! a\,\mathrm{Li}_2(a)
       + 6\log a\,\mathrm{Li}_3(a) - 6\,\mathrm{Li}_4(a).
\]
The route above never differentiates $\mathrm{Li}_4$.  Substituting
$u = t/2$ \emph{first} collapses $W_0$, since $\log^2(t/2)$ becomes
$\log^2 u$; the remaining factor is then an exact differential in four of
the six rows,
\[
\frac{H_1}{t}\,dt = d\bigl[\mathrm{Li}_2(t)\bigr],\quad
\frac{H_1}{1-t}\,dt = d\Bigl[\tfrac{H_1^2}{2}\Bigr],\quad
\frac{H_2}{t}\,dt = d\bigl[\mathrm{Li}_2(u)\bigr],\quad
\frac{H_2}{2-t}\,dt = d\Bigl[\tfrac{H_2^2}{2}\Bigr],
\]
so the dilogarithm cross terms integrate by the chain rule,
$\int \mathrm{Li}_2\,d\mathrm{Li}_2 = \mathrm{Li}_2^2/2$, with no series
at all; and where a series is needed, $\mathrm{Li}_4(1/2)$ arrives as
$\sum 2^{-n}n^{-4}$, that is, as the \emph{answer} rather than as
something an antiderivative must produce.  The two relations
\eqref{eq:R12}--\eqref{eq:R21} then avoid $\mathrm{Li}_4(1/2)$ entirely.

One tempting shortcut is a dead end and is recorded to save the reader the
trouble.  The bracket in \eqref{eq:six} \emph{is} an exact differential:
its primitive is the M\"obius pullback of the integrand of
\eqref{eq:coeff}.  Integrating \eqref{eq:six} by parts against
$W_0' = -2R/t$ therefore reverses the substitution that produced it and
returns \eqref{eq:coeff}.  Nothing is evaluated.

\subsection{Verification status}

Formulas \eqref{eq:coeff}, \eqref{eq:six}, \eqref{eq:K}, the reduction of
$K$ to the two cubic-linear sums, and the rows $I_{1t}$, $I_{1,1-t}$,
$I_{2t}$, $I_{2,2-t}$ are proved in Lean~4 (toolchain \texttt{v4.29.0})
and depend only on \texttt{propext}, \texttt{Classical.choice} and
\texttt{Quot.sound}.  The relations \eqref{eq:R12}--\eqref{eq:R21}, hence
the two remaining rows, are at present verified exactly over $\mathbb{Q}$
against the four proved rows and numerically to fifteen digits, but are
not yet formalized.  The six row values were additionally obtained by
three mutually independent analytic derivations and cross-checked against
high-precision quadrature.

We note the scope precisely: \eqref{eq:S} is one of the seven scalar
evaluations that Stage 2 requires.  Three of the others are likewise
proved in the formalization; the remaining three are not.
\bibitem{Schneider2007}
C.~Schneider,
\emph{Symbolic summation assists combinatorics},
S\'em.\ Lotharingien Combin.\ \textbf{56} (2007), B56b.

\bibitem{BBG2001}
J.~M.~Borwein, D.~M.~Bradley, and D.~J.~Broadhurst,
\emph{Evaluation of $k$-fold Euler/Zagier sums},
Electron.\ J.\ Combin.\ \textbf{4} (2001), R5.
\end{thebibliography}

\end{document}
